\chapter*{Notation and Conventions}
\addcontentsline{toc}{chapter}{Notation and Conventions}
\markboth{Notation and Conventions}{Notation and Conventions}
\begingroup
\small
\setlength{\abovedisplayskip}{6pt}
\setlength{\belowdisplayskip}{6pt}
\setlength{\abovedisplayshortskip}{3pt}
\setlength{\belowdisplayshortskip}{4pt}
\renewcommand{\paragraph}[1]{\par\addvspace{6pt}\noindent\textbf{#1}\enspace}
Most notation is introduced again when it first becomes mathematically
relevant. The conventions collected here are those used repeatedly
throughout the notes, especially where closely related meanings coexist.

\paragraph{Numbers and maps.}
We use $\N=\{1,2,3,\ldots\}$ and $\N_0=\N\cup\{0\}$.
The symbols $\Z,\Q,\R,\C$ denote the integers, rationals, reals, and
complex numbers. General maps usually use $f,g,h$; linear maps and
operators use $T,S,L$. Charts and coordinate maps use $\varphi,\psi,\chi$,
curves use $\gamma$, and surface parametrizations usually use $\sigma$.
Conventional symbols such as a force field $F$ are retained.
Composition is $(g\circ f)(x)=g(f(x))$; $f|_A$ is the
restriction to $A$.

\paragraph{Euclidean inputs and scalar volume.}
In Part II the full Euclidean input is written
\[
\mathbf{x}=(x_1,\ldots,x_n),\qquad f(\mathbf{x}),\qquad
\int_Qf(\mathbf{x})\,\dd\mathbf{x}.
\]
Components $x_i$, scalar coordinates, and abstract points retain ordinary
type. The unsigned scalar volume symbol $\dd\mathbf{x}$ is distinct from
coordinate covectors $\dd x_i$, their oriented wedge product, and exterior
differentiation. Established vector and manifold notation is retained.

\paragraph{Four related differential notations.}
\[
\begin{array}{c|l}
Df(a)&\text{Euclidean/Fr\'echet derivative at }a\\
(df)_p&\text{value of the exact scalar-function $1$-form at }p\\
df_p:T_pM\to T_{f(p)}N&\text{intrinsic differential of a smooth map }f:M\to N\\
d\omega&\text{exterior derivative of a differential form}
\end{array}
\]
For a Euclidean scalar function, $(df)_p=Df(p)$ under the usual
tangent-space identification. For a manifold scalar function, identify
$T_{f(p)}\R$ with $\R$ to regard its differential as a covector.
\[
\boxed{\begin{gathered}
D\text{ emphasizes Euclidean differentiation;}\\
d\text{ is used for intrinsic differentials and forms.}
\end{gathered}}
\]
\paragraph{Bases and their measurements.}
A basis $e_1,\ldots,e_n$ has dual basis $e^1,\ldots,e^n$, with
$e^i(e_j)=\delta^i_j$. The Kronecker delta $\delta^i_j$ is $1$ when
$i=j$ and $0$ otherwise. In manifold coordinates the corresponding pair is
\[
\left.\frac{\partial}{\partial x^i}\right|_p,
\qquad dx^i|_p.
\]
The first is a tangent direction; the second measures its $i$th coordinate.
A topological basis, by contrast, is a family of open sets, not vectors.

\paragraph{Sets and boundaries.}
The symbols $\overline A$, $A^\circ$, and $\partial A$ denote closure,
interior, and topological boundary in the stated ambient space, with
$\partial A=\overline A\setminus A^\circ$. The intrinsic manifold boundary
$\partial M$ is defined by half-space charts and need not equal the
boundary of $M$ as a subset of an ambient Euclidean space. For example,
a circle has empty intrinsic boundary, while its topological boundary in
$\R^2$ is the entire circle. Context and the specified ambient space
matter. The notation $\operatorname{supp}f$ or $\operatorname{supp}\omega$
means the closure of the set where the function or form is nonzero.
\par
\endgroup
